% ==================================================
% CHAPTER 1: PROJECT OVERVIEW AND CONTEXT
% ==================================================

\chapter{Project Overview and Context}
\label{ch:overview}

\section{Software Bid and Selection}
\label{sec:bid}

The project was proposed in August 2026 as a \emph{personal web memory}: a
browser companion that keeps the text of pages a person reads and lets them
find a page again by describing it. The proposal, discussed at the first
meeting with the supervisor on 17~August, set out the flow that the final
system still follows: capture a page in the browser, embed and index its
text, and answer a plain-language query with the matching page.

\begin{table}[H]
\centering
\small
\begin{tabularx}{\textwidth}{>{\raggedright\arraybackslash}p{0.22\textwidth} L}
\toprule
\textbf{Bid item} & \textbf{Content} \\
\midrule
Problem & People cannot find pages they have read: browser history and bookmarks match only titles and addresses. \\
Proposed product & A Chrome extension, a web app and a backend that store page text and search it by meaning and by keywords. \\
Team & Rabeeh Ahmad and Ishan Singhania (group 3C51), supervised by Dr.\ Raghav B.\ Venkataramaiyer. \\
First milestone & A working backend that returns the right URLs within three days of approval (met on 19~August). \\
Effort and schedule & One semester (July--November 2026), delivered in iterations reported weekly on the course forum. \\
Resources & Open-source software; one cloud server (AWS EC2 \texttt{c7i-flex.large}); free tiers of Amazon Cognito, Amazon SES and Google Gemini. \\
\bottomrule
\end{tabularx}
\caption{Summary of the software bid}
\label{tab:bid}
\end{table}

\paragraph{Selection rationale.} The project was chosen because the problem
is common among students like us, its value is easy to demonstrate, and it
exercises the whole software-engineering
cycle: requirements with privacy and security constraints, a design with
several cooperating parts, an implementation that can be deployed and used,
and search quality that can be \emph{measured} with standard metrics rather
than judged by eye.

\section{Problem, Purpose and Objectives}

\paragraph{Problem.} Students, developers and researchers read hundreds of
web pages a week, and much of it is useful again later. By then they usually
remember the idea, not the title, the website or the exact words. Browser
history records \emph{that} a page was visited and can be searched only by
title and address; bookmarks must be made in advance; web search engines
search the whole internet rather than the pages a person has read.

\paragraph{Purpose.} MemoryLane treats the pages a person reads as a
searchable personal memory: it stores their readable text and lets the user
find a page by describing it in their own words, for example
\textit{``the article about reducing database query latency''} or
\textit{``the reddit post on rust lifetimes from last week''}.

\paragraph{Formal statement.} Let user $u$ own a library
$L_u = \{p_1, \dots, p_n\}$ of saved pages. Given a natural-language
description $q$ of a target page $p^{\star} \in L_u$, return a ranking
$R(q) \subseteq L_u$ with $|R(q)| \le 10$ that maximises the mean reciprocal
rank of the target,
\[
  \operatorname{E}\left[\frac{1}{\operatorname{rank}\bigl(p^{\star} \mid R(q)\bigr)}\right]
  \qquad \text{(MRR@10; the term is 0 if } p^{\star} \notin R(q)\text{)},
\]
subject to:
\begin{enumerate}[nosep]
  \item \textbf{Isolation:} $\forall p \in R(q):\ \operatorname{owner}(p) = u$.
  \item \textbf{Placement:} each page lives whole in one store
  $s(p) \in \{\text{device}, \text{server}\}$, and
  $\operatorname{score}(q, p \mid \text{device}) = \operatorname{score}(q, p \mid \text{server})$.
  \item \textbf{Privacy:} pages on the device and backups are encrypted; the
  server holds no key to a backup.
  \item \textbf{Latency:} median response time $\le 100$\,ms.
\end{enumerate}
Each part is measured: MRR@10 and latency on SciFact
(Section~\ref{sec:evaluation}) and placement by the parity checks
(Section~\ref{sec:local-engine}).

\paragraph{Objectives.}
\begin{enumerate}
  \item O1: Save pages manually and, optionally, automatically, without ever
  storing private pages.
  \item O2: Clean, chunk and embed saved pages so they can be searched by
  meaning and by exact words.
  \item O3: Return the right page first for most descriptive queries, with
  readable snippets and optional summaries.
  \item O4: Keep every user's library private, and deploy the system
  publicly over HTTPS.
  \item O5: Measure retrieval quality with standard metrics and use the
  results to improve the system.
  \item O6: Follow sound engineering practice: specified requirements,
  recorded design decisions, automated tests and documentation.
\end{enumerate}

\paragraph{Audience.} This report is written for the course evaluators and
supervisor, and for developers who would maintain or extend MemoryLane.

\section{Scope and System Boundary}

\paragraph{In scope.} Manual and automatic saving of web pages and YouTube
videos from Chrome; cleaning, chunking, embedding and indexing; hybrid search
with site and time hints, snippets and on-demand summaries; a library with
deletion and undo; accounts with Google or email sign-in; problem reports;
keeping pages on the device (encrypted, searched in the browser) with the
server as a fallback, and encrypted backups; deployment on one server;
documentation and evaluation.

\paragraph{Out of scope.} Crawling the web; sharing pages between users;
browsers other than Chromium-based ones; mobile apps; importing existing
history or bookmarks; pages behind logins that the user's browser cannot
read; offline use.

\paragraph{Boundary.} The system consists of the extension (with its local
engine and IndexedDB store), the web app and the server (Caddy, the FastAPI
backend, Solr and three SQLite databases). The
reader's browser, Amazon Cognito, Google Gemini, Amazon SES and YouTube are
outside it and are reached only through defined interfaces
(Section~\ref{sec:interfaces}); Figure~\ref{fig:dfd0} shows the boundary as a
context diagram.

\section{Product Perspective}

MemoryLane is a new, self-contained product that complements, rather than
replaces, existing tools:

\begin{table}[H]
\centering
\small
\begin{tabularx}{\textwidth}{>{\raggedright\arraybackslash}p{0.2\textwidth} L L}
\toprule
\textbf{Existing system} & \textbf{What it does} & \textbf{What MemoryLane adds} \\
\midrule
Browser history & Lists visited addresses and titles & Search by the page's content and meaning \\
Bookmarks & Saves chosen pages in folders & Saving without organising; automatic capture \\
Web search engines & Search the whole web & Search only the pages you have read \\
Read-later services & Keep pages to read later & Find pages already read, by description \\
\bottomrule
\end{tabularx}
\caption{MemoryLane compared with existing systems}
\label{tab:existing}
\end{table}

Its major functions are saving, searching, summarising and managing pages.
It depends on Amazon Cognito for identity, on Google Gemini for summaries
(optional) and on Amazon SES for report email (optional), and runs on
Apache Solr, FastAPI and Caddy (Figure~\ref{fig:components}).

\section{Stakeholders and Users}

\begin{table}[H]
\centering
\small
\begin{tabularx}{\textwidth}{>{\raggedright\arraybackslash}p{0.15\textwidth} L L L L}
\toprule
\textbf{Class} & \textbf{Role} & \textbf{Needs} & \textbf{Skills} & \textbf{Access} \\
\midrule
Reader & Saves and finds pages (primary user) & Fast, private search; no effort to save & Ordinary browser use & Own library only, after sign-in \\
Maintainer & Deploys and operates the service & Health status, problem reports, safe updates & Docker, Linux, AWS & Server, configuration, reports (not users' pages through the app) \\
Supervisor and evaluators & Assess the project & Clear requirements, evidence, results & Software engineering & Reports, documentation, repository \\
External services & Identity, summaries, email & Well-formed requests & -- & Only what each interface sends \\
\bottomrule
\end{tabularx}
\caption{Stakeholders and user classes}
\label{tab:stakeholders}
\end{table}

\section{Operating Environment}

\begin{table}[H]
\centering
\small
\begin{tabularx}{\textwidth}{>{\raggedright\arraybackslash}p{0.17\textwidth} L}
\toprule
\textbf{Part} & \textbf{Environment} \\
\midrule
Extension & Google Chrome or another Chromium browser with Manifest V3 support, on any desktop OS; WebGPU used when present for saving on the device (WebAssembly otherwise); ~130\,MB model cached once \\
Web app & Any modern browser on desktop or phone; voice search where the browser supports speech recognition (Chrome, Edge) \\
Server & Linux with Docker Compose; AWS EC2 \texttt{c7i-flex.large} (2 vCPU, 4\,GB RAM, 2\,GB swap), Mumbai region; no GPU \\
Network & HTTPS only (ports 80/443 open); ordinary home or campus connections; the server reaches Cognito, Gemini and SES over the internet \\
\bottomrule
\end{tabularx}
\caption{Operating environment}
\label{tab:environment}
\end{table}

\section{Constraints}

\begin{itemize}
  \item \textbf{Technology}: one server without a GPU, so embedding runs on
  the CPU in one process; Solr 9 for both keyword and vector search; a Chrome
  extension (Manifest V3), installed unpacked rather than from a store.
  \item \textbf{Resources}: a two-person team over one semester, and
  free-tier cloud services.
  \item \textbf{Security policy}: no client secret in the browser; every
  request authenticated; no user may reach another's pages.
  \item \textbf{Manifest V3}: no code may be loaded at run time, so the
  extension bundles Pyodide and its packages and the model's runtime; only
  the model's weights (data) are downloaded.
  \item \textbf{Privacy and legal}: page text is personal data, so by default
  it stays on the device; on the server it is kept out of logs, deletable at any time, and
  sent to an external model only when the user asks for a summary. The
  SciFact evaluation data is licensed CC BY-NC 2.0 and is used for evaluation
  only. YouTube's restrictions on server requests mean captions are fetched in
  the user's browser.
\end{itemize}

\section{Assumptions and Dependencies}

\begin{table}[H]
\centering
\small
\begin{tabularx}{\textwidth}{L L}
\toprule
\textbf{Assumption or dependency} & \textbf{Effect if it fails} \\
\midrule
Users sign in before saving & Nothing is saved; the extension shows a sign-in screen \\
Amazon Cognito is reachable & Known sign-ins keep working; a first-time sign-in gets \texttt{503} rather than a wrong library \\
Google Gemini is configured and reachable & Summaries are unavailable; search and everything else work \\
Amazon SES is reachable & Problem reports wait in SQLite and are retried with back-off \\
YouTube serves captions to the user's browser & A video is saved from its description, or skipped \\
Pages expose readable text & Pages with under 50 words are refused with a clear message \\
The server has its Docker volumes & Without them, pages and sign-in links would be lost; EBS snapshots are the backup \\
The browser keeps its IndexedDB data & Device pages are lost if the data is cleared; the encrypted backup restores them \\
\bottomrule
\end{tabularx}
\caption{Assumptions, dependencies and the effect of their failure}
\label{tab:assumptions}
\end{table}

\section{Terminology and Report Guide}

\paragraph{Terminology.} Abbreviations are listed before this chapter. A
\emph{page} is a saved web page or video; a \emph{chunk} is a passage of
about 150 words; a \emph{library} is all the pages of one person. An
\emph{embedding} is a list of 384 numbers representing what a chunk means.
\emph{BM25} ranks chunks by the query's words, weighting rare words more;
\emph{hybrid search} runs BM25 and embedding (KNN) search together, and
\emph{RRF} combines their ranked lists by adding $1/(60+\text{rank})$ from
each. \emph{HNSW} is the graph index Solr uses to find the nearest embeddings
quickly. The \emph{relevance floor} drops results that are neither a strong
keyword match nor close enough in meaning.

\paragraph{Notation.} Requirements and tests carry identifiers used
throughout the report: \texttt{O} objectives, \texttt{FR} functional
requirements, \texttt{QR} quality requirements, \texttt{BR} business rules,
\texttt{UC} use cases, \texttt{US} user stories and \texttt{TC} test cases.
ADR numbers refer to the architecture decision records.

\paragraph{Report guide.} Chapter~\ref{ch:requirements} specifies the
requirements and models them with use cases, stories, activity diagrams and
data-flow diagrams. Chapter~\ref{ch:design} presents the design: architecture
and decisions, class, sequence, collaboration and state chart diagrams.
Chapter~\ref{ch:implementation} describes the implementation and delivery:
components, deployment, interfaces, storage, screenshots, build and release,
documentation and the changes since the prototype.
Chapter~\ref{ch:testing} covers the test plan, test cases, results, defects
and coverage, and Chapter~\ref{ch:results} the evaluation results, challenges,
limitations, individual contributions and conclusions.
